% © 2026 Ing. David Jaroš — CC BY-NC-ND 4.0
%
% File: research_tracks/T3_ALPHA/theta0_electroweak_vacuum_orbit_theorem.tex
% Purpose: Close the remaining alpha proof condition by showing that any
%          non-zero UBT electroweak vacuum in the minimal Theta doublet sector
%          is gauge-equivalent to the standard (SU(2)_L,Y=1) doublet vacuum.
% Status: [L1 conditional] — closes the representation/orbit part of EW-2.
%         Remaining assumption: the symmetry-breaking potential V(Theta) has a
%         non-zero minimum in the minimal electroweak Theta doublet sector.
% Date: 2026-06-13

\ifdefined\INCLUDEMODE\else
\documentclass[12pt,a4paper]{article}
\usepackage[a4paper,margin=2.5cm]{geometry}
\usepackage{amsmath,amssymb,amsthm,mathtools}
\usepackage{xcolor}
\usepackage[most]{tcolorbox}
\usepackage{hyperref}

\newtheorem{definition}{Definition}[section]
\newtheorem{lemma}[definition]{Lemma}
\newtheorem{theorem}[definition]{Theorem}
\newtheorem{corollary}[definition]{Corollary}
\newtheorem{remark}[definition]{Remark}

\newtcolorbox{statusbox}[1][]{
  colback=yellow!8!white,
  colframe=orange!70!black,
  arc=3pt,
  boxrule=1pt,
  title=Status,
  #1
}
\newtcolorbox{resultbox}[1][]{
  colback=green!8!white,
  colframe=green!50!black,
  arc=3pt,
  boxrule=1pt,
  title=Result,
  #1
}

\newcommand{\Lone}{\textbf{[L1 cond.]}}

\title{The $\Theta_0$ Electroweak Vacuum Orbit Theorem\\
       \large Closing the Representation Gap in the UBT Alpha Track}
\author{Ing.\ David Jaro\v{s}}
\date{2026-06-13}

\begin{document}
\maketitle
\fi

\begin{abstract}
The alpha proof requires that the UBT electromagnetic projection be the compact
unit-charge \(U(1)_{\rm EM}\) line bundle.  The previous EM projection theorem
reduced this to the condition that the low-energy UBT vacuum transform as an
\((SU(2)_L,Y=1)\) electroweak doublet.  This note proves the representation and
orbit part of that condition.  Since
\(\mathbb C\otimes\mathbb H\cong {\rm Mat}(2,\mathbb C)\), the minimal
left-action electroweak component of \(\Theta\) is a complex \(SU(2)_L\)
doublet.  The right scalar phase gives primitive hypercharge weight \(Y=1\).
For any non-zero minimum in this sector, \(SU(2)_L\) acts transitively on the
norm sphere, so the vacuum is gauge-equivalent to
\[
  \Phi_0=\begin{pmatrix}0\\ v/\sqrt2\end{pmatrix},
  \qquad Y\Phi_0=\Phi_0.
\]
Thus the standard electroweak doublet vacuum is not an extra shape assumption;
it is the canonical gauge representative of any non-zero minimal UBT doublet
vacuum.  The only remaining physical input is the existence of a non-zero
symmetry-breaking minimum of \(V(\Theta)\) in this sector.
\end{abstract}

\begin{statusbox}
\textbf{Classification.} This closes the representation/orbit part of the
previous EW-2 condition.  It does not derive the scale \(v\) or prove from first
principles that \(V(\Theta)\) must have a non-zero electroweak minimum.  It shows
that if such a minimum exists in the minimal UBT electroweak sector, its doublet
form and hypercharge are forced by the algebra and gauge orbit.
\end{statusbox}

\section{Minimal electroweak sector of $\Theta$}

The biquaternion algebra has the standard matrix representation
\[
  \mathbb C\otimes\mathbb H \cong {\rm Mat}(2,\mathbb C).
\]
The left unitary action gives \(SU(2)_L\).  A minimal electroweak component of
\(\Theta\) may therefore be represented by a complex column
\[
  \Phi=\begin{pmatrix}\phi_1\\ \phi_2\end{pmatrix}\in\mathbb C^2,
\]
with transformation law
\[
  \Phi\mapsto U_L\Phi,
  \qquad U_L\in SU(2)_L.
\]
The right scalar phase acts as
\[
  \Phi\mapsto e^{i\beta}\Phi.
\]
With the standard electroweak convention
\[
  \Phi\mapsto e^{i\beta Y/2}\Phi,
\]
this is the primitive weight \(Y=1\) doublet.

\section{Orbit theorem}

\begin{lemma}[Transitivity of $SU(2)$ on the non-zero doublet sphere]
For every non-zero \(\Phi\in\mathbb C^2\) there exists \(U\in SU(2)\) such that
\[
  U\Phi=\begin{pmatrix}0\\ \|\Phi\|\end{pmatrix}.
\]
\end{lemma}

\begin{proof}
Write \(\Phi=(a,b)^T\) and \(r=\sqrt{|a|^2+|b|^2}\).  For \(r>0\), the matrix
\[
  U=\frac1r
  \begin{pmatrix}
    b & -a\\
    \bar a & \bar b
  \end{pmatrix}
\]
is unitary and has determinant one, hence \(U\in SU(2)\).  Direct
multiplication gives \(U\Phi=(0,r)^T\).  Therefore all non-zero doublets of the
same norm lie in one \(SU(2)\) orbit.
\end{proof}

\begin{theorem}[$\Theta_0$ electroweak vacuum orbit \Lone]
Assume the UBT symmetry-breaking potential \(V(\Theta)\) has a non-zero minimum
in the minimal electroweak doublet sector.  Then a gauge representative of the
vacuum is
\[
  \Phi_0=\begin{pmatrix}0\\ v/\sqrt2\end{pmatrix},
  \qquad Y\Phi_0=1\cdot\Phi_0.
\]
Thus the UBT low-energy vacuum has the standard \((SU(2)_L,Y=1)\) electroweak
doublet form required by the EM projection theorem.
\end{theorem}

\begin{proof}
Let \(\Phi_*\ne0\) be a vacuum minimiser in the minimal electroweak sector.
Gauge invariance of \(V\) implies that all elements of the \(SU(2)_L\) orbit of
\(\Phi_*\) have the same vacuum energy.  By the transitivity lemma, there exists
\(U\in SU(2)_L\) such that
\[
  U\Phi_*=\begin{pmatrix}0\\ \|\Phi_*\|\end{pmatrix}.
\]
Writing \(\|\Phi_*\|=v/\sqrt2\) gives the standard representative
\[
  \Phi_0=\begin{pmatrix}0\\v/\sqrt2\end{pmatrix}.
\]
The right scalar phase action fixes the primitive hypercharge weight \(Y=1\),
so \(Y\Phi_0=\Phi_0\).
\end{proof}

\section{Consequence for the alpha proof}

The previous EM projection theorem showed that for
\[
  \Phi_0=\begin{pmatrix}0\\v/\sqrt2\end{pmatrix},
  \qquad Y=1,
\]
the unique unbroken generator is
\[
  Q_{\rm EM}=T_3+\frac{Y}{2}.
\]
This note proves that this vacuum shape is the canonical gauge representative of
any non-zero UBT doublet vacuum.  Therefore the alpha proof no longer needs to
postulate the shape of \(\Theta_0\).  It only needs the standard symmetry-breaking
condition:
\[
  \boxed{V(\Theta)\text{ has a non-zero minimum in the minimal electroweak doublet sector.}}
\]

\begin{resultbox}
\textbf{Updated alpha status.} The final chain is now:
\[
  V(\Theta)\text{ non-zero EW doublet minimum}
  \Longrightarrow
  \Theta_0\sim(0,v/\sqrt2),\;Y=1
  \Longrightarrow
  Q_{\rm EM}=T_3+Y/2
  \Longrightarrow
  U(1)_{\rm EM}\text{ compact unit charge}
  \Longrightarrow
  n=\alpha^{-1}.
\]
The remaining gap is not the EM projection or charge normalisation; it is the
existence of the non-zero electroweak minimum from \(V(\Theta)\).
\end{resultbox}

\ifdefined\INCLUDEMODE\else
\end{document}
\fi
